\begin{proposition}[Conjugate members and antiholomorphic automorphisms]\hypertarget{tag-01QG}{}\label{fam:src:B.4.2}\label{fam:st:B.4.2}
\StatementRecord{fam:src:B.4.2}{proved}{01QG}{\SeriesRef{fam:src:B.4.1}, \SeriesRef{fam:src:A.3.1}, \SeriesRef{fam:src:B.2.1}, \SeriesRef{fam:src:A.5.1}}

Use the normalization $(\beta^{\rm part}+\tau)(p_0)=0$. There is an antiholomorphic diffeomorphism $X_c\to X_{c'}$ exactly when
\[
 c'-(1-\bar c)\in2\mathbb Z,
 \qquad\text{equivalently }w'=-\bar w.
\]
Consequently $\overline{X_c}\simeq X_{1-\bar c}$. The members admitting antiholomorphic self-maps are exactly those with $\operatorname{Re}c\in\tfrac12+\mathbb Z$, equivalently $w=-\bar w$. On each such member the group of holomorphic and antiholomorphic diffeomorphisms is
\[
 \operatorname{Aut}^{\pm}(X_c)\simeq\mathbb C^*\rtimes\mathbb Z/2,
 \qquad z\longmapsto1/\bar z
\]
for the action of the second factor. Its antiholomorphic maps form the nonidentity coset. For all other members the full group is the holomorphic automorphism group of Theorem~\ref{fam:src:B.4.1}.
\end{proposition}


\begin{proof}
\begin{enumerate}[label=\arabic*.,leftmargin=*]
\item We first construct the comparison before returning its affine data to the canonical choice. Lift conjugation of the real orbifold base to the antiholomorphic involution $\iota$ of its universal cover whose cusp branch has $\tau(\iota z)=-\overline{\tau(z)}$. A lift exists by the universal covering property; it extends at the finite orbifold points, whose local equations are $s_j^{m_j}=t_j$. Choose the usual paths from a real point $t>1$, so reflection fixes the second centre and sends the first chosen lift to $g_2^{-1}z_1$. The induced involution of the deck group is
\[
 \phi(g_1)=g_2^{-1}g_1^{-1}g_2,\qquad
 \phi(g_2)=g_2^{-1},\qquad \phi(g_0)=g_0^{-1}.
\]
These identities also follow by reflecting the lifted peripheral paths; their product and their squares agree with $g_1g_2g_0=1$. The cusp branch fixes the lift uniquely.

\item Denote the affine period transformations by $G_1,G_2$ on the triple $(\tau,\mu,\beta)$ and put $K(\tau,\mu,\beta)=(-\bar\tau,1-\bar\mu,-\bar\beta)$. Direct rational substitution gives
\[
 KG_1K=G_2^{-1}G_1^{-1}G_2,\qquad KG_2K=G_2^{-1}.
\]
For example $G_2^{-1}(\tau,\mu,\beta)=(-1/\tau,(1-\mu)/\tau,\beta+3-6(1-\mu)^2/\tau)$ and
$G_1^{-1}(\tau,\mu,\beta)=(1/(1-\tau),1-\mu/(1-\tau),\beta-2+6\mu^2/(1-\tau))$. These formulas verify every constant in the two equalities. The reflected $\mu$ is cusp bounded and has the same affine laws, so their difference is a global section of $\mathcal O_B(-p_0)$ and vanishes. After this equality, the difference between the reflected and original $\beta^{\rm part}$ is a holomorphic invariant base function, hence constant; its cusp value is zero under the stated normalization. Therefore
\[
 \mu(\iota z)=1-\bar\mu(z),\qquad
 \beta^{\rm part}(\iota z)=-\overline{\beta^{\rm part}(z)}.
\]
The original $\tau$ identity follows from its $j$ equation and its specified cusp branch. In particular $D_{-\bar c}(\iota z)=D_c(z)$, so this comparison stays in the exact admissible domain.

\item Let
\begin{equation}
 H=\begin{pmatrix}1&0&0&0\\0&1&0&0\\-6&0&-1&0\\0&-1&0&-1\end{pmatrix}.
 \label{I:eq:real-reflection-matrix}
\end{equation}
Then $H^2=I$, $HA_2H=A_2^{-1}$, $HA_1H=A_2^{-1}A_1^{-1}A_2$, and
\[
 \Pi_{-\bar c}(\iota z)H=-\overline{\Pi_c(z)}.
\]
Thus $(z,\zeta)\mapsto(\iota z,-\bar\zeta)$ is a well-defined antiholomorphic torus map. It also descends under the middle deck action. Indeed covariance and the displayed lattice conjugacies imply $R_{\phi(g)}(\iota z)=\overline{R_g(z)}$ by multiplying the full real-invertible period matrices. Thus the reflection intertwines the middle deck actions.

\item At the first finite centre compose this map with the target deck identification $g_2$ to return the centre $g_2^{-1}z_1$ to $z_1$. Its lattice map becomes $A_2H$ and it conjugates the clockwise generator to the inverse generator. At the second centre its lattice map is $H$ with the same reversal. Consequently the target affine vectors are exactly
\begin{equation}
 w'_1=-A_2H\varepsilon=-\varepsilon+d,\qquad
 w'_2=-H(-\varepsilon')=\varepsilon'-3d.
 \label{I:eq:real-reflected-affine-data}
\end{equation}
The minus signs arise because the inverse of a generator with invariant translation $w'_j/m_j$ has translation $-w'_j/m_j$. This calculation verifies the full affine actions on the root covers. For the overlaps choose the root coordinates of the real local base coordinates $t$ and $t-1$ compatible with these local reflections. This choice has an explicit comparison with any original root coordinate: their ratio is a holomorphic unit, so changing it changes the collar by $(2\pi i)^{-1}\log h\,\Pi w_j$, a translation extending on the entire root family. Its inverse gives a holomorphic isomorphism of the two glued descriptions. Thus the extension property and final classification transfer to the original coordinates. Then $\log s'_j=\overline{\log s_j}$ on paired branches. Under $\zeta\mapsto-\bar\zeta$, conjugation of $1/(2\pi i)$ supplies the compensating sign, and the identity of the two logarithmic collar translations follows exactly. Two roots of the same fixed real coordinate differ by an $m_j$th root of unity; this changes the comparison only by a constant period-coordinate translation, and the additive correction of \ref{fam:src:A.3.1} glues those constants. The preceding holomorphic-unit comparison handles a different original local base coordinate. Its inverse has the same construction.

\item At the cusp choose $s=(2\pi i)^{-1}\log t_c$, so $s\mapsto-\bar s$. The full holomorphic unit matrix satisfies
\[
 C_{-\bar c}(\bar t_c)=-\overline{C_c(t_c)}+
       \begin{pmatrix}6&0\\0&1\end{pmatrix}.
\]
The integral extra matrix has trivial exponential. Hence $(x,t_c)\mapsto(\bar x,\bar t_c)$ intertwines the full cusp deck groups, with the same integral index $\lambda$. The toric fan is defined over $\mathbb Z$, so this map extends antiholomorphically on every toric chart and stratum. It preserves the untwisted cusp collar. We have therefore constructed an antiholomorphic map to the target datum with parameter $-\bar c$, $r'_1=-1,k'_1=1,r'_2=1,k'_2=-3,l'_0=0$, whose $p'$ is $+1$.

\item Apply \ref{fam:src:B.2.1} to normalize this target back to the canonical $p=-1$ datum. Take $a=-1$ and the odd shift $n=1$. This gives a holomorphic map to $X_{1-\bar c}$. Every other antiholomorphic map, composed with the inverse of this one, is a holomorphic comparison between canonical members. \ref{fam:src:A.5.1} gives exactly the stated congruence. Exponentiating it gives $w'=-\bar w$.

\item  For a real member, choose any resulting antiholomorphic self-map $J$. Its derivative on ordinary fibre coordinates is $+\text{complex conjugation}$: the initial reflection has derivative $-\text{complex conjugation}$, while the normalization just used has linear part $-I_2$. Hence it conjugates translation by $b\in\mathbb C/\mathbb Z$ to translation by $\bar b$. Its square is a holomorphic translation $b_0$. Commuting $J$ with $J^2$ shows $b_0=\bar b_0$ in $\mathbb C/\mathbb Z$. This equality forces the imaginary part of a representative of $b_0$ to vanish, so take $b_0\in\mathbb R$. Replacing $J$ by the composition of translation by $-b_0/2$ with $J$ changes its square to $b_0-b_0/2-b_0/2=0$. We have an antiholomorphic involution. Under $b\mapsto e^{2\pi ib}$, its conjugation action is $z\mapsto1/\bar z$. This proves the split group and specifies its antiholomorphic coset.
\end{enumerate}
\end{proof}
